\documentclass[a4paper,12pt]{article}

\usepackage[T2A]{fontenc}
\usepackage[utf8]{inputenc}
\usepackage[russian]{babel}

\usepackage{amsmath,amssymb,amsthm,mathtools}
\usepackage{helvet}
\renewcommand{\familydefault}{\sfdefault}

\usepackage{icomma}
\usepackage[a4paper,left=20mm,right=20mm,top=20mm,bottom=22mm]{geometry}
\usepackage{microtype}

\usepackage{tikz}
\usetikzlibrary{calc}
\usepackage{tkz-euclide}

\usepackage{pgfplots}
\pgfplotsset{compat=1.18}

\usepackage{tabularx,booktabs,longtable}
\usepackage{enumitem}
\usepackage{hyperref}
\usepackage{parskip}
\usepackage{fancyhdr}
\usepackage{setspace}
\usepackage{xcolor}

% ------------------------------------------------------------
% Палитра
% ------------------------------------------------------------
\definecolor{accent}{HTML}{008080}
\definecolor{darkaccent}{HTML}{004D4D}
\definecolor{textgray}{HTML}{333333}
\definecolor{softgray}{HTML}{707070}
\definecolor{linegray}{HTML}{D5DDDD}

% ------------------------------------------------------------
% Настройки
% ------------------------------------------------------------
\geometry{headheight=15pt}
\onehalfspacing
\setlength{\parindent}{0pt}
\setlength{\parskip}{0.55em}

\hypersetup{
    colorlinks=true,
    linkcolor=darkaccent,
    urlcolor=darkaccent,
    pdftitle={Чек-лист. Тригонометрия и алгебраический минимум для ОГЭ},
    pdfauthor={Лёвин Артём Александрович}
}

\pagestyle{fancy}
\fancyhf{}
\fancyhead[L]{\small\textcolor{softgray}{ОГЭ по математике}}
\fancyhead[R]{\small\textcolor{softgray}{София Хоменко}}
\fancyfoot[C]{\small\textcolor{softgray}{\thepage}}
\renewcommand{\headrulewidth}{0.3pt}
\renewcommand{\headrule}{\hbox to\headwidth{\color{linegray}\leaders\hrule height \headrulewidth\hfill}}

\setlist[itemize]{
    leftmargin=1.5em,
    itemsep=0.25em,
    topsep=0.25em
}

\setlist[enumerate]{
    leftmargin=1.8em,
    itemsep=0.45em,
    topsep=0.3em
}

\DeclareMathOperator{\tg}{tg}
\DeclareMathOperator{\ctg}{ctg}

\newcommand{\checkitem}{\item[$\square$]}
\newcommand{\accentline}{%
    \vspace{0.2em}
    {\color{accent}\rule{\linewidth}{0.6pt}}
    \vspace{0.5em}
}

\newcommand{\topic}[1]{%
    \vspace{0.5em}
    {\large\bfseries\color{darkaccent}#1}\par
    \accentline
}

\begin{document}

% ============================================================
% ТИТУЛЬНЫЙ ЛИСТ
% ============================================================
\begin{titlepage}
\thispagestyle{empty}

\begin{center}

\vspace*{2.2cm}

{\color{accent}\rule{0.18\linewidth}{1.2pt}}

\vspace{1.1cm}

{\Huge\bfseries\color{darkaccent}
Чек-лист}

\vspace{0.45cm}

{\LARGE
Тригонометрия и алгебраический минимум\\[0.2cm]
для ОГЭ}

\vspace{1.4cm}

{\large
Табличные значения тригонометрических функций\\
Формулы приведения\\
Прямоугольный треугольник\\
Базовые алгебраические преобразования
}

\vspace{2.1cm}

{\large\bfseries
Лёвин Артём Александрович}

\vspace{0.3cm}

{\large
эксклюзивно для Софии Хоменко}

\vspace{1.3cm}

{\large \today}

\vfill

\begin{tikzpicture}[x=1cm,y=1cm]
    \draw[
        accent,
        line width=1.1pt
    ]
    (-7.5,0)
        .. controls (-4.7,0.65) and (-2.7,-0.45) ..
    (0,0)
        .. controls (2.8,0.45) and (5.0,-0.55) ..
    (7.5,0.05);
\end{tikzpicture}

\vspace{0.35cm}

{\small\color{softgray}
Краткое повторение ключевых идей занятия
}

\end{center}
\end{titlepage}

% ============================================================
% ОСНОВНОЙ ЧЕК-ЛИСТ
% ============================================================

\section*{\color{darkaccent}1. Табличные значения синуса и косинуса}

\topic{Что нужно знать наизусть}

Для углов
\[
0^\circ,\quad 30^\circ,\quad 45^\circ,\quad 60^\circ,\quad 90^\circ
\]
значения синуса и косинуса являются базовыми.

\begin{center}
\renewcommand{\arraystretch}{1.4}
\begin{tabular}{c c c c c c}
\toprule
$\alpha$ &
$0^\circ$ &
$30^\circ$ &
$45^\circ$ &
$60^\circ$ &
$90^\circ$
\\
\midrule
$\sin\alpha$ &
$0$ &
$\dfrac12$ &
$\dfrac{\sqrt2}{2}$ &
$\dfrac{\sqrt3}{2}$ &
$1$
\\[0.25cm]
$\cos\alpha$ &
$1$ &
$\dfrac{\sqrt3}{2}$ &
$\dfrac{\sqrt2}{2}$ &
$\dfrac12$ &
$0$
\\
\bottomrule
\end{tabular}
\end{center}

\topic{Способ быстро восстановить таблицу}

Для синуса удобно помнить последовательность
\[
\sin\alpha=
\frac{\sqrt{0}}2,\quad
\frac{\sqrt{1}}2,\quad
\frac{\sqrt{2}}2,\quad
\frac{\sqrt{3}}2,\quad
\frac{\sqrt{4}}2.
\]

Она соответствует углам
\[
0^\circ,\ 30^\circ,\ 45^\circ,\ 60^\circ,\ 90^\circ.
\]

Для косинуса порядок обратный:
\[
\cos\alpha=
\frac{\sqrt{4}}2,\quad
\frac{\sqrt{3}}2,\quad
\frac{\sqrt{2}}2,\quad
\frac{\sqrt{1}}2,\quad
\frac{\sqrt{0}}2.
\]

Например,
\[
\sin60^\circ=\frac{\sqrt3}{2},
\qquad
\cos30^\circ=\frac{\sqrt3}{2}.
\]

\begin{itemize}
\checkitem Я знаю таблицу для $0^\circ$, $30^\circ$, $45^\circ$, $60^\circ$, $90^\circ$.
\checkitem Я могу восстановить забытое значение через формулу $\dfrac{\sqrt n}{2}$.
\checkitem Я помню, что $\sin60^\circ=\cos30^\circ$.
\checkitem Я помню, что $\sin30^\circ=\cos60^\circ$.
\end{itemize}

% ------------------------------------------------------------

\section*{\color{darkaccent}2. Формулы приведения}

\topic{Главные формулы с $180^\circ$}

Пусть $\alpha$ --- острый угол.

\[
\boxed{\sin(180^\circ-\alpha)=\sin\alpha}
\]

\[
\boxed{\cos(180^\circ-\alpha)=-\cos\alpha}
\]

\[
\boxed{\tg(180^\circ-\alpha)=-\tg\alpha}
\]

\[
\boxed{\ctg(180^\circ-\alpha)=-\ctg\alpha}
\]

Значения $\tg\alpha$ и $\ctg\alpha$ рассматриваются при тех углах,
где соответствующая функция определена.

\topic{Как использовать формулу}

Если угол отсутствует в таблице, удобно представить его через
$180^\circ-\alpha$.

Например,
\[
120^\circ=180^\circ-60^\circ.
\]

Следовательно,
\[
\sin120^\circ
=
\sin(180^\circ-60^\circ)
=
\sin60^\circ
=
\frac{\sqrt3}{2}.
\]

Для косинуса знак меняется:
\[
\cos150^\circ
=
\cos(180^\circ-30^\circ)
=
-\cos30^\circ
=
-\frac{\sqrt3}{2}.
\]

Ещё один пример:
\[
\sin135^\circ
=
\sin(180^\circ-45^\circ)
=
\sin45^\circ
=
\frac{\sqrt2}{2}.
\]

\topic{Короткая памятка}

\begin{center}
\renewcommand{\arraystretch}{1.45}
\begin{tabularx}{0.84\linewidth}{>{\raggedright\arraybackslash}X c}
\toprule
Выражение & Результат\\
\midrule
$\sin(180^\circ-\alpha)$ & $\sin\alpha$\\
$\cos(180^\circ-\alpha)$ & $-\cos\alpha$\\
$\tg(180^\circ-\alpha)$ & $-\tg\alpha$\\
$\ctg(180^\circ-\alpha)$ & $-\ctg\alpha$\\
\bottomrule
\end{tabularx}
\end{center}

Для запоминания полезно заметить:

\[
\sin(180^\circ-\alpha)
\]
сохраняет положительный знак, а у косинуса, тангенса и котангенса
появляется минус.

\topic{Формулы с $90^\circ$}

Полезно узнавать и такие формы:
\[
\sin(90^\circ+\alpha)=\cos\alpha,
\]
\[
\cos(90^\circ+\alpha)=-\sin\alpha.
\]

Например,
\[
\sin135^\circ
=
\sin(90^\circ+45^\circ)
=
\cos45^\circ
=
\frac{\sqrt2}{2}.
\]

Для текущего этапа подготовки основной акцент делаем на формулах с
$180^\circ$.

\begin{itemize}
\checkitem Я умею представить $120^\circ$ как $180^\circ-60^\circ$.
\checkitem Я умею представить $135^\circ$ как $180^\circ-45^\circ$.
\checkitem Я умею представить $150^\circ$ как $180^\circ-30^\circ$.
\checkitem Я контролирую знак перед косинусом, тангенсом и котангенсом.
\end{itemize}

% ============================================================

\section*{\color{darkaccent}3. Тригонометрия в прямоугольном треугольнике}

\topic{Определения}

Рассмотрим прямоугольный треугольник $ABC$ с прямым углом $B$.

\begin{center}
\begin{tikzpicture}[scale=1.05]
    \tkzSetUpLine[line width=1pt,color=accent]
    \tkzSetUpPoint[color=accent,fill=accent,size=3]
    \tkzSetUpLabel[color=black,font=\small]

    \tkzDefPoint(0,0){A}
    \tkzDefPoint(4.6,0){B}
    \tkzDefPoint(4.6,3.2){C}

    \tkzDrawPolygon(A,B,C)
    \tkzDrawPoints(A,B,C)

    \tkzMarkRightAngle[size=0.32](A,B,C)
    \tkzMarkAngle[size=0.65](B,A,C)
    \tkzLabelAngle[pos=0.95](B,A,C){$\alpha$}

    \tkzLabelPoints[below](A,B)
    \tkzLabelPoints[right](C)
\end{tikzpicture}
\end{center}

Для угла $\alpha=\angle BAC$:

\[
\boxed{
\sin\alpha=
\frac{\text{противолежащий катет}}
     {\text{гипотенуза}}
}
\]

\[
\boxed{
\cos\alpha=
\frac{\text{прилежащий катет}}
     {\text{гипотенуза}}
}
\]

\[
\boxed{
\tg\alpha=
\frac{\text{противолежащий катет}}
     {\text{прилежащий катет}}
}
\]

На данном чертеже:
\[
\sin\angle BAC=\frac{BC}{AC},
\]
\[
\cos\angle BAC=\frac{AB}{AC},
\]
\[
\tg\angle BAC=\frac{BC}{AB}.
\]

\begin{itemize}
\checkitem Я сначала определяю гипотенузу.
\checkitem Я определяю катет, лежащий напротив выбранного угла.
\checkitem Я определяю катет, прилежащий к выбранному углу.
\checkitem Я записываю нужное отношение сторон.
\end{itemize}

% ------------------------------------------------------------

\section*{\color{darkaccent}4. Внешний угол и формулы приведения}

\topic{Главная конструкция}

Пусть луч $AF$ является продолжением стороны $AB$.

\begin{center}
\begin{tikzpicture}[scale=1.0]
    \tkzSetUpLine[line width=1pt,color=accent]
    \tkzSetUpPoint[color=accent,fill=accent,size=3]
    \tkzSetUpLabel[color=black,font=\small]

    \tkzDefPoint(-2.2,0){F}
    \tkzDefPoint(0,0){A}
    \tkzDefPoint(4.8,0){B}
    \tkzDefPoint(4.8,3.3){C}

    \tkzDrawSegments(F,B A,C B,C)
    \tkzDrawPoints(F,A,B,C)

    \tkzMarkRightAngle[size=0.32](A,B,C)

    \tkzMarkAngle[size=0.62](B,A,C)
    \tkzLabelAngle[pos=0.95](B,A,C){$\alpha$}

    \tkzMarkAngle[size=0.83](C,A,F)
    \tkzLabelAngle[pos=1.15](C,A,F){$\varphi$}

    \tkzLabelPoints[below](F,A,B)
    \tkzLabelPoints[right](C)
\end{tikzpicture}
\end{center}

Углы $\alpha$ и $\varphi$ смежные, поэтому
\[
\varphi=180^\circ-\alpha.
\]

Тогда
\[
\sin\varphi
=
\sin(180^\circ-\alpha)
=
\sin\alpha,
\]
а
\[
\cos\varphi
=
\cos(180^\circ-\alpha)
=
-\cos\alpha.
\]

Следовательно,
\[
\boxed{
\sin\varphi=\frac{BC}{AC}
}
\]
и
\[
\boxed{
\cos\varphi=-\frac{AB}{AC}
}.
\]

\topic{Алгоритм}

Если требуется тригонометрическая функция внешнего угла:

\begin{enumerate}
    \item найдите связанный с ним внутренний угол;
    \item запишите
    \[
    \varphi=180^\circ-\alpha;
    \]
    \item примените формулу приведения;
    \item перейдите к отношению сторон прямоугольного треугольника;
    \item при необходимости найдите недостающую сторону по теореме Пифагора.
\end{enumerate}

% ============================================================

\section*{\color{darkaccent}5. Теорема Пифагора как первый инструмент}

\topic{Что проверять при виде прямоугольного треугольника}

Если катеты равны $a$ и $b$, а гипотенуза равна $c$, то
\[
\boxed{a^2+b^2=c^2}.
\]

Например, при
\[
c=5,\qquad a=4
\]
получаем
\[
b^2=5^2-4^2=25-16=9,
\]
откуда
\[
b=3.
\]

После этого можно вычислить:
\[
\sin\alpha,\qquad
\cos\alpha,\qquad
\tg\alpha.
\]

\begin{itemize}
\checkitem При виде прямоугольного треугольника я сразу проверяю возможность применения теоремы Пифагора.
\checkitem Я помню, что длина стороны положительна.
\checkitem Я умею по двум сторонам восстановить третью.
\end{itemize}

% ============================================================

\section*{\color{darkaccent}6. Равнобедренный треугольник и высота}

\topic{Полезное свойство}

В равнобедренном треугольнике высота, проведённая из вершины к основанию,
одновременно является медианой.

Пусть $ABC$ --- равнобедренный треугольник,
\[
AB=BC,
\]
а $BH$ --- высота к основанию $AC$.

Тогда
\[
AH=HC=\frac{AC}{2}.
\]

\begin{center}
\begin{tikzpicture}[scale=1.0]
    \tkzSetUpLine[line width=1pt,color=accent]
    \tkzSetUpPoint[color=accent,fill=accent,size=3]
    \tkzSetUpLabel[color=black,font=\small]

    \tkzDefPoint(-3.5,0){A}
    \tkzDefPoint(0,3.8){B}
    \tkzDefPoint(3.5,0){C}
    \tkzDefPoint(0,0){H}

    \tkzDrawPolygon(A,B,C)
    \tkzDrawSegment(B,H)
    \tkzDrawPoints(A,B,C,H)

    \tkzMarkRightAngle[size=0.32](A,H,B)

    \tkzLabelPoints[below](A,H,C)
    \tkzLabelPoints[above](B)
\end{tikzpicture}
\end{center}

После проведения высоты появляются два прямоугольных треугольника.
В каждом из них можно использовать:

\[
a^2+b^2=c^2,
\]
\[
\sin\alpha,
\qquad
\cos\alpha,
\qquad
\tg\alpha.
\]

\topic{Пример}

Пусть
\[
BC=5,\qquad AC=8.
\]

Тогда
\[
HC=\frac{AC}{2}=4.
\]

Из прямоугольного треугольника $BHC$:
\[
BH^2+HC^2=BC^2,
\]
\[
BH^2+16=25,
\]
\[
BH=3.
\]

Следовательно, для угла при основании:
\[
\tg\angle BCA
=
\frac{BH}{HC}
=
\frac34.
\]

% ============================================================

\section*{\color{darkaccent}7. Квадратные уравнения: краткое повторение}

\topic{Через дискриминант}

Для
\[
ax^2+bx+c=0,\qquad a\ne0
\]
используем
\[
D=b^2-4ac.
\]

Если $D>0$, то
\[
x_{1,2}
=
\frac{-b\pm\sqrt D}{2a}.
\]

Например:
\[
x^2-x-6=0.
\]

Здесь
\[
a=1,\qquad b=-1,\qquad c=-6,
\]
поэтому
\[
D=(-1)^2-4\cdot1\cdot(-6)=25.
\]

Тогда
\[
x_1=\frac{1+5}{2}=3,
\qquad
x_2=\frac{1-5}{2}=-2.
\]

Ответ:
\[
x=-2,\quad x=3.
\]

\topic{Через вынесение общего множителя}

\[
x^2-x=0,
\]
\[
x(x-1)=0.
\]

По свойству произведения:
\[
x=0
\quad\text{или}\quad
x=1.
\]

\topic{Разность квадратов}

\[
x^2-6=0
\]
можно записать как
\[
x^2-(\sqrt6)^2=0.
\]

Получаем
\[
(x-\sqrt6)(x+\sqrt6)=0,
\]
откуда
\[
x=\pm\sqrt6.
\]

Для уравнения
\[
x^2+6=0
\]
действительных решений нет, поскольку
\[
x^2\ge0
\]
для любого действительного $x$.

% ============================================================

\section*{\color{darkaccent}8. Разложение квадратного трёхчлена}

Если $x_1$ и $x_2$ --- корни квадратного трёхчлена
\[
ax^2+bx+c,
\]
то
\[
\boxed{
ax^2+bx+c
=
a(x-x_1)(x-x_2)
}.
\]

Например, корни выражения
\[
x^2-x-6
\]
равны $3$ и $-2$.

Значит,
\[
x^2-x-6
=
(x-3)(x+2).
\]

\begin{itemize}
\checkitem Я сначала нахожу корни.
\checkitem Я учитываю коэффициент $a$ перед $x^2$.
\checkitem При отрицательном корне внимательно работаю с двойным минусом.
\end{itemize}

% ============================================================

\section*{\color{darkaccent}9. Формулы сокращённого умножения}

\topic{Разность квадратов}

\[
\boxed{a^2-b^2=(a-b)(a+b)}
\]

Например,
\[
9y^2-16
=
(3y)^2-4^2
=
(3y-4)(3y+4).
\]

\topic{Квадрат разности}

\[
\boxed{(a-b)^2=a^2-2ab+b^2}
\]

Например,
\[
x^2-4x+4=(x-2)^2.
\]

Полезно проверять три элемента:

\[
a^2,\qquad -2ab,\qquad b^2.
\]

% ============================================================

\section*{\color{darkaccent}10. Алгебраические дроби и знак «минус»}

\topic{ОДЗ}

Перед сокращением дроби обязательно фиксируем условие:
\[
\text{знаменатель}\ne0.
\]

Например,
\[
\frac{3x+2}{-3x-2}.
\]

ОДЗ:
\[
-3x-2\ne0,
\]
\[
x\ne-\frac23.
\]

Знаменатель можно представить как
\[
-3x-2=-(3x+2).
\]

Поэтому
\[
\frac{3x+2}{-3x-2}
=
\frac{3x+2}{-(3x+2)}
=
-1,
\qquad
x\ne-\frac23.
\]

\topic{«Странствующий» минус}

При $B\ne0$:
\[
\boxed{
\frac{A}{-B}
=
-\frac{A}{B}
=
\frac{-A}{B}
}.
\]

Если минус переносится в числитель, он относится ко всему числителю:
\[
-\frac{4x+2}{5x+1}
=
\frac{-(4x+2)}{5x+1}.
\]

Скобки сохраняют правильный смысл выражения.

\topic{Одинаковые знаменатели через вынесение минуса}

Рассмотрим:
\[
\frac{1}{4x-2}
-
\frac{1}{2-4x}.
\]

ОДЗ:
\[
4x-2\ne0,
\]
то есть
\[
x\ne\frac12.
\]

Заметим:
\[
2-4x=-(4x-2).
\]

Тогда
\[
-\frac{1}{2-4x}
=
-\frac{1}{-(4x-2)}
=
\frac{1}{4x-2}.
\]

Следовательно,
\[
\frac{1}{4x-2}
-
\frac{1}{2-4x}
=
\frac{1}{4x-2}
+
\frac{1}{4x-2}
=
\frac{2}{4x-2}.
\]

Сокращаем:
\[
\boxed{
\frac{1}{2x-1}
},
\qquad
x\ne\frac12.
\]

\begin{itemize}
\checkitem Перед работой с дробью я записываю ОДЗ.
\checkitem Я умею вынести минус из знаменателя.
\checkitem Я понимаю, к какому выражению относится знак минус.
\checkitem Сокращаю только множители.
\checkitem После преобразований сохраняю исходные ограничения.
\end{itemize}

% ============================================================

\section*{\color{darkaccent}11. Итоговый алгоритм}

\begin{enumerate}
    \item Если встречается табличный угол, вспоминаю значение сразу.
    \item Если угол больше $90^\circ$, проверяю возможность представить его как
    \[
    180^\circ-\alpha.
    \]
    \item При внешнем угле связываю его со внутренним через сумму $180^\circ$.
    \item В прямоугольном треугольнике определяю гипотенузу и нужные катеты.
    \item Проверяю, требуется ли теорема Пифагора.
    \item В равнобедренном треугольнике использую свойства высоты к основанию.
    \item В алгебре сначала ищу структуру:
    общий множитель, разность квадратов, квадратный трёхчлен.
    \item В алгебраических дробях начинаю с ОДЗ.
    \item Все знаки «минус» контролирую отдельным шагом.
\end{enumerate}

% ============================================================
% ТРЕНИРОВОЧНЫЙ БЛОК
% ============================================================

\clearpage
\section*{\color{darkaccent}Тренировочный блок}

\begin{center}
{\large\bfseries Путь А: 5 упражнений}
\end{center}

\accentline

\textbf{Упражнение 1. Табличные значения и формулы приведения}

Вычислите:

\begin{enumerate}[label=\alph*)]
    \item $\sin30^\circ$;
    \item $\cos60^\circ$;
    \item $\sin135^\circ$;
    \item $\cos120^\circ$;
    \item $\cos150^\circ$.
\end{enumerate}

\vspace{0.7em}

\textbf{Упражнение 2. Внешний угол прямоугольного треугольника}

В прямоугольном треугольнике $ABC$
\[
\angle B=90^\circ,\qquad
AB=8,\qquad
BC=6.
\]

Луч $AF$ является продолжением луча $AB$ в противоположную сторону.

Найдите:
\[
\sin\angle FAC
\qquad\text{и}\qquad
\cos\angle FAC.
\]

В решении обязательно укажите длину гипотенузы $AC$.

\vspace{0.7em}

\textbf{Упражнение 3. Теорема Пифагора и тригонометрия}

В прямоугольном треугольнике $KLM$
\[
\angle L=90^\circ,\qquad
KM=13,\qquad
KL=5.
\]

Найдите:
\[
LM,
\qquad
\sin\angle K,
\qquad
\cos\angle K,
\qquad
\tg\angle K.
\]

\vspace{0.7em}

\textbf{Упражнение 4. Равнобедренный треугольник}

В равнобедренном треугольнике $ABC$
\[
AB=BC=10,\qquad AC=12.
\]

Высота $BH$ проведена к основанию $AC$.

Найдите:

\begin{enumerate}[label=\alph*)]
    \item $AH$;
    \item $BH$;
    \item $\sin\angle BAC$;
    \item $\cos\angle BAC$;
    \item $\tg\angle BAC$.
\end{enumerate}

\vspace{0.7em}

\textbf{Упражнение 5. Алгебраический блок}

Выполните задания.

\begin{enumerate}[label=\alph*)]
    \item Решите уравнение
    \[
    x^2-5x-14=0.
    \]

    \item Разложите на множители:
    \[
    25y^2-9.
    \]

    \item Разложите на множители:
    \[
    x^2-5x-14.
    \]

    \item Упростите выражение и укажите ОДЗ:
    \[
    \frac{5x-3}{3-5x}.
    \]

    \item Упростите выражение и укажите ОДЗ:
    \[
    \frac{1}{6x-4}
    -
    \frac{1}{4-6x}.
    \]
\end{enumerate}

% ============================================================
% ОТВЕТЫ
% ============================================================

\clearpage
\section*{\color{darkaccent}Ответы к тренировочному блоку}

\renewcommand{\arraystretch}{1.5}

\begin{table}[ht]
\centering
\caption{Краткая таблица ответов}
\begin{tabularx}{\linewidth}{
    >{\raggedright\arraybackslash}p{0.12\linewidth}
    >{\raggedright\arraybackslash}X
}
\toprule
\textbf{№} & \textbf{Ответ}\\
\midrule

1 &
\[
\text{а) }\frac12;
\qquad
\text{б) }\frac12;
\qquad
\text{в) }\frac{\sqrt2}{2};
\qquad
\text{г) }-\frac12;
\qquad
\text{д) }-\frac{\sqrt3}{2}.
\]
\\

\midrule

2 &
\[
AC=10,
\qquad
\sin\angle FAC=\frac35,
\qquad
\cos\angle FAC=-\frac45.
\]
\\

\midrule

3 &
\[
LM=12,
\qquad
\sin\angle K=\frac{12}{13},
\qquad
\cos\angle K=\frac5{13},
\qquad
\tg\angle K=\frac{12}{5}.
\]
\\

\midrule

4 &
\[
AH=6,
\qquad
BH=8,
\]
\[
\sin\angle BAC=\frac45,
\qquad
\cos\angle BAC=\frac35,
\qquad
\tg\angle BAC=\frac43.
\]
\\

\midrule

5 &
\text{а)}
\[
x=-2,\quad x=7.
\]

\text{б)}
\[
25y^2-9=(5y-3)(5y+3).
\]

\text{в)}
\[
x^2-5x-14=(x-7)(x+2).
\]

\text{г)}
\[
-1,
\qquad
x\ne\frac35.
\]

\text{д)}
\[
\frac{1}{3x-2},
\qquad
x\ne\frac23.
\]
\\

\bottomrule
\end{tabularx}
\end{table}

\vfill

\begin{center}
\begin{tikzpicture}[x=1cm,y=1cm]
    \draw[
        accent,
        line width=0.9pt
    ]
    (-6.7,0)
        .. controls (-4.0,0.45) and (-2.1,-0.3) ..
    (0,0)
        .. controls (2.4,0.3) and (4.2,-0.45) ..
    (6.7,0);
\end{tikzpicture}

\vspace{0.35cm}

{\small\color{softgray}
Лёвин Артём Александрович
}
\end{center}

\end{document}